% Beamer Presentation Basics
% Demonstrates: slides, title page, bullet points, columns, and blocks.
% Compile with: pdflatex main.tex

\documentclass[aspectratio=169]{beamer}

% --- Theme and colors ---
\usetheme{Madrid}
\usecolortheme{default}

% --- Metadata ---
\title{Introduction to Beamer}
\subtitle{Creating Presentations with \LaTeX{}}
\author{Your Name}
\institute{Your University}
\date{\today}

\begin{document}

% --- Title slide ---
\begin{frame}
\titlepage
\end{frame}

% --- Table of contents ---
\begin{frame}{Outline}
\tableofcontents
\end{frame}

% --- Section 1 ---
\section{Getting Started}

\begin{frame}{What is Beamer?}
\begin{itemize}
    \item Beamer is a \LaTeX{} document class for creating presentations
    \item Produces PDF slides with professional formatting
    \item Supports themes, animations, and overlays
    \item Used extensively in academia and research
\end{itemize}
\end{frame}

% --- Section 2: Columns and blocks ---
\section{Layout Features}

\begin{frame}{Two-Column Layout}
\begin{columns}
    \begin{column}{0.5\textwidth}
        \textbf{Left Column}
        \begin{itemize}
            \item Use \texttt{columns} environment
            \item Specify width as fraction of \texttt{\textbackslash textwidth}
        \end{itemize}
    \end{column}
    \begin{column}{0.5\textwidth}
        \textbf{Right Column}
        \begin{enumerate}
            \item Numbered lists work too
            \item Combine with other environments
        \end{enumerate}
    \end{column}
\end{columns}
\end{frame}

\begin{frame}{Blocks}
\begin{block}{Standard Block}
    This is a standard block with a blue header.
\end{block}

\begin{alertblock}{Alert Block}
    This is an alert block --- use it for warnings or key points.
\end{alertblock}

\begin{exampleblock}{Example Block}
    This is an example block --- use it for examples or code.
\end{exampleblock}
\end{frame}

% --- Section 3: Math ---
\section{Mathematics}

\begin{frame}{Equations in Slides}
Beamer supports all standard \LaTeX{} math:

\begin{equation}
    e^{i\pi} + 1 = 0
\end{equation}

The quadratic formula:
\begin{equation}
    x = \frac{-b \pm \sqrt{b^2 - 4ac}}{2a}
\end{equation}
\end{frame}

% --- Final slide ---
\begin{frame}{}
\centering
\Huge Thank you!

\vspace{1cm}
\normalsize Questions?
\end{frame}

\end{document}
